\chapter{Wybór technologii}
\label{rozdzial3}

Wybór odpowiednich technologii jest kluczowym etapem procesu wytwarzania oprogramowania, mającym bezpośredni wpływ na wydajność, skalowalność, bezpieczeństwo oraz łatwość późniejszego utrzymania systemu (maintainability). Decyzje podjęte na tym etapie determinują szybkość procesu deweloperskiego oraz komfort pracy programistów. W niniejszym rozdziale przedstawiono szczegółową charakterystykę stosu technologicznego wykorzystanego do realizacji projektu ,,Household Manager''.

\section{Backend — Go, GORM, PostgreSQL, Chi Router}

Warstwa serwerowa (backend) stanowi mózg aplikacji, odpowiadający za logikę biznesową, przetwarzanie danych oraz komunikację z bazą danych. Do jej implementacji wybrano język \textbf{Go} (znany również jako Golang).

\subsection{Uzasadnienie wyboru Go}

Wybór języka Go (Golang) jako fundamentu warstwy serwerowej podyktowany był potrzebą zapewnienia wysokiej wydajności oraz niezawodności przy jednoczesnej prostocie implementacji logiki biznesowej. Jako język nowoczesny, rozwijany przez firmę Google, Go adresuje problemy znane ze starszych technologii, takie jak długi czas kompilacji czy skomplikowane zarządzanie pamięcią.


Główne powody wyboru tej technologii to:

\begin{enumerate}
    \item \textbf{Wydajność kompilowana:} Go jest językiem kompilowanym bezpośrednio do kodu maszynowego, co sprawia, że aplikacje napisane w Go są znacznie szybsze niż te uruchamiane w środowiskach interpretowanych (jak Python czy JavaScript/Node.js). Nie wymaga maszyny wirtualnej (JVM), co przekłada się na mniejsze zużycie pamięci RAM i szybszy start aplikacji.
    
    \item \textbf{Statyczne typowanie:} Silny system typów pomaga wykrywać błędy już na etapie kompilacji, zanim kod trafi na produkcję. W przeciwieństwie do języków dynamicznych, programista ma pewność, że zmienna zadeklarowana jako liczba całkowita zawsze nią będzie, co zwiększa stabilność systemu.
    
    \item \textbf{Współbieżność (Concurrency):} Go posiada unikalny model współbieżności oparty na tzw. \textit{goroutines} i kanałach (\textit{channels}). Goroutine to "lekki wątek", zarządzany przez runtime języka, a nie przez system operacyjny. Uruchomienie tysięcy goroutines wymaga ułamka zasobów potrzebnych dla tradycyjnych wątków systemowych. W kontekście serwera HTTP oznacza to możliwość jednoczesnej obsługi ogromnej liczby zapytań od użytkowników przy minimalnym narzucie sprzętowym.
    
    \item \textbf{Prostota i czytelność:} Składnia Go jest minimalistyczna (brak klas, dziedziczenia w tradycyjnym sensie, skomplikowanych meta-programowania). Wymusza to pisanie prostego, czytelnego kodu, który jest łatwy do zrozumienia dla nowych członków zespołu. Sformatowany kod narzędziem \texttt{gofmt} wygląda zawsze tak samo, co eliminuje spory o styl kodowania.
\end{enumerate}

W porównaniu do popularnego Node.js, Go oferuje lepszą wydajność obliczeniową (istotne np. przy przetwarzaniu obrazów w module OCR) oraz stabilniejszy ekosystem bibliotek standardowych.

\subsection{Framework i biblioteki backendowe}

W ekosystemie Go panuje tendencja do unikania rozbudowanych frameworków na rzecz komponowania aplikacji z mniejszych, wyspecjalizowanych bibliotek. W projekcie wykorzystano następujące rozwiązania:

\begin{itemize}
    \item \textbf{Chi Router:}
    Lekki i szybki router HTTP, w pełni kompatybilny ze standardową biblioteką \texttt{net/http}. Wykorzystano go do budowy przejrzystego systemu \textit{middleware}, odpowiadającego za logowanie żądań, obsługę polityki CORS oraz weryfikację tokenów JWT. Chi zapewnia wysoką czytelność kodu przy minimalnym narzucie wydajnościowym.
    
    \item \textbf{GORM (Go Object Relational Mapping):}
    Biblioteka typu ORM, która umożliwiła mapowanie struktur języka Go na relacje w bazie PostgreSQL. GORM został wykorzystany do automatyzacji operacji CRUD, obsługi złożonych relacji między encjami (np. \textit{many-to-many} dla członków domu) oraz zarządzania migracjami schematu bazy danych.
    
    \item \textbf{Google Gemini API (Generative AI):}
    Zamiast klasycznych silników OCR, w systemie wykorzystano zaawansowany model multimodalny \textit{Gemini 1.5 Flash}. Pozwala on nie tylko na odczyt tekstu (OCR), ale również na jego semantyczną analizę i strukturyzację danych w formacie JSON.
    
    
    
    \item \textbf{Google Cloud SDK / HTTP Client:}
    Do komunikacji z modelem Gemini wykorzystano natywny klient HTTP języka Go. Implementacja (usługa \texttt{OCRService}) obejmuje bezpieczne przesyłanie obrazów zakodowanych w standardzie Base64 oraz autorski algorytm \texttt{calculateGeminiConfidence}, który ocenia wiarygodność zwróconych danych (np. poprzez weryfikację sumy pozycji paragonu z kwotą całkowitą).
\end{itemize}
W pracy dyplomowej najlepiej wygląda układ, w którym każda kluczowa technologia ma swój dedykowany podpunkt. Pozwala to recenzentowi szybko odnaleźć konkretne narzędzia i ocenić, czy student rozumie ich rolę w projekcie.

Oto kompletna, rozbudowana sekcja o froncie, przygotowana w sposób najbardziej "dyplomowy" (oddzielne podpunkty dla każdego narzędzia).

\section{Frontend — React Native, Expo, TypeScript}

Warstwa kliencka (aplikacja mobilna) została zrealizowana w architekturze opartej na komponentach, co zapewnia modularność kodu i łatwość rozbudowy interfejsu użytkownika.

\subsection{React Native i architektura komponentowa}

React Native, stworzony przez firmę Meta, jest frameworkiem pozwalającym na budowanie natywnych aplikacji mobilnych przy użyciu biblioteki React. W przeciwieństwie do rozwiązań hybrydowych, React Native nie korzysta z widoków przeglądarkowych (\textit{WebView}), lecz mapuje komponenty JavaScript na natywne widoki systemów iOS i Android. Zapewnia to wysoką wydajność interfejsu oraz dostęp do płynnych animacji systemowych. Wykorzystanie paradygmatu deklaratywnego UI pozwala na automatyczne aktualizowanie widoku w odpowiedzi na zmiany stanu aplikacji.

\subsection{Język TypeScript}

W projekcie zdecydowano się na użycie języka \textbf{TypeScript} jako nadzbioru języka JavaScript. TypeScript wprowadza statyczne typowanie, co w kontekście pracy inżynierskiej przynosi następujące korzyści:
\begin{itemize}
\item \textbf{Wykrywanie błędów w czasie kompilacji:} Eliminuje typowe błędy typu \textit{runtime}, takie jak próba dostępu do pola w obiekcie o wartości \textit{null}.
\item \textbf{Dokumentacja poprzez kod:} Interfejsy i typy jasno definiują strukturę danych przesyłanych między ekranami aplikacji.
\item \textbf{Refaktoryzacja:} Dzięki silnemu typowaniu, wprowadzanie zmian w dużych strukturach danych (np. modelu użytkownika) jest bezpieczniejsze i szybsze.
\end{itemize}

\subsection{Platforma Expo i usługi EAS}

Projekt oparto na środowisku \textbf{Expo} – ekosystemie narzędzi ułatwiających proces wytwarzania aplikacji w React Native. Expo dostarcza ujednolicone API do funkcji systemowych (aparat, galeria zdjęć, system plików). Wykorzystano również \textbf{EAS (Expo Application Services)}, co umożliwiło generowanie plików instalacyjnych (\textit{.apk} dla systemu Android) w chmurze, bez konieczności posiadania wydajnej stacji roboczej skonfigurowanej pod natywny proces kompilacji.

\subsection{Stylizacja interfejsu — NativeWind}

Do budowy warstwy wizualnej wykorzystano bibliotekę \textbf{NativeWind}, która implementuje system klas narzędziowych \textbf{Tailwind CSS} w środowisku mobilnym. Podejście to (\textit{utility-first}) pozwala na szybkie stylowanie komponentów bezpośrednio w kodzie JSX, co drastycznie skraca czas tworzenia responsywnych widoków i zapewnia spójność wizualną całego systemu bez konieczności pisania zewnętrznych arkuszy stylów.

\subsection{Nawigacja — React Navigation}

Zarządzanie przepływem ekranów zrealizowano przy pomocy biblioteki \textbf{React Navigation}. W systemie zaimplementowano dwa rodzaje nawigatorów:
\begin{itemize}
\item \textbf{Stack Navigator:} Wykorzystywany do liniowych procesów, takich jak logowanie, rejestracja czy przejście do szczegółów konkretnego rachunku.
\item \textbf{Bottom Tab Navigator:} Stanowi główne menu aplikacji, umożliwiając szybkie przełączanie się między modułami: Dom, Zadania, Finanse i Profil.
\end{itemize}

\subsection{Zarządzanie stanem globalnym — Zustand}

Do przechowywania danych dostępnych w całej aplikacji (np. informacja o zalogowanym użytkowniku, identyfikator wybranego domu) wybrano bibliotekę \textbf{Zustand}. Jest to rozwiązanie minimalistyczne, oparte na \textit{hookach}, które oferuje wysoką wydajność przy znacznie mniejszej złożoności niż tradycyjne biblioteki typu Redux. Pozwala to na uniknięcie problemu \textit{prop drilling} (przekazywania danych przez wiele poziomów komponentów).

\subsection{Komunikacja sieciowa — Axios}

Interakcja z backendem realizowana jest za pomocą klienta \textbf{Axios}. Biblioteka ta wspiera mechanizm \textit{interceptors}, który w projekcie został wykorzystany do automatycznego wstrzykiwania tokena JWT do nagłówków każdego żądania HTTP oraz do globalnej obsługi błędów sieciowych (np. automatyczne wylogowanie w przypadku wygaśnięcia sesji).

\subsection{Bezpieczeństwo — Expo Secure Store}

Ze względu na wrażliwy charakter danych dostępowych, w projekcie zrezygnowano z domyślnego \textit{AsyncStorage} na rzecz \textbf{Expo Secure Store}. Biblioteka ta zapewnia szyfrowane przechowywanie par klucz-wartość na poziomie systemowym (Keychain na iOS, SharedPreferences z szyfrowaniem na Androidzie), co gwarantuje, że tokeny autoryzacyjne użytkownika są chronione przed niepowołanym dostępem innych aplikacji.

\section{Integracja backend-frontend (REST API)}

Komunikacja między aplikacją mobilną a serwerem odbywa się zgodnie z zasadami architektury \textbf{REST} (Representational State Transfer).
Główne założenia REST spełnione w projekcie:
\begin{itemize}
    \item \textbf{Bezstanowość (Stateless):} Każde żądanie HTTP zawiera wszystkie informacje niezbędne do jego obsłużenia (np. token autoryzacyjny). Serwer nie przechowuje stanu sesji użytkownika między zapytaniami.
    \item \textbf{Jednolity interfejs:} Zasoby są identyfikowane przez adresy URL (np. \texttt{/api/homes}), a operacje na nich wykonywane są za pomocą standardowych metod HTTP (GET, POST, PUT, DELETE).
    \item \textbf{Format danych:} Wymiana danych odbywa się w lekkim formacie JSON, który jest natywnie wspierany zarówno przez JavaScript (Frontend), jak i biblioteki Go (Backend).
\end{itemize}

Wybrano REST zamiast GraphQL ze względu na prostotę implementacji, łatwość cachowania zapytań HTTP oraz jasny podział odpowiedzialności.

\section{Testy jednostkowe i integracyjne w Go}

Jakość kodu serwerowego jest weryfikowana za pomocą testów. Wykorzystano bibliotekę \textbf{Testify}, która rozszerza standardowy pakiet \texttt{testing} o asercje (np. \texttt{assert.Equal}, \texttt{assert.Nil}) oraz narzędzia do mockowania zależności. Testy obejmują kluczowe elementy logiki biznesowej oraz handlery HTTP.

\section{Podsumowanie}

Dobór technologii w projekcie ,,Household Manager'' podyktowany był chęcią stworzenia nowoczesnego, skalowalnego i łatwego w utrzymaniu systemu. Połączenie wydajnego backendu w Go z elastycznym frontendem w React Native stanowi solidną podstawę dla aplikacji użytkowej.
